\subsection{Module D.18: CE1 one-gap reverse-path certificate}
\label{subsec:app-ce1-direct-certificate}

\paragraph{Exact interface.}
Use the exact seven-point witness and the five center-free handoffs of
\zcref{prop:readable-k410-ce1}.  Only the retained radial reaches at
$T_4,T_3,T_2$ enter; no radial point on $r_0$ is part of this interface.

\paragraph{Variables and feasible set.}
Take the signed CE1 variables
$0<R<1$, $w=1-R$, $E=\sqrt{1-Rw}$, $\eta=1-E$, and $P=E(1-E)$.
The positive slacks $A=\alpha$ and $D=\delta$ satisfy the two signed CE1
inequalities displayed below.  The five V roles have the stated type,
handoff, and three closed radial-demand constraints.

\paragraph{Objective.}
Starting from a hypothetical upper bound on the returning actual reach
$A_0$, propagate the selected high-radial output twice and force a final
contradiction at $T_1$.

\paragraph{Exact result.}

\begin{proposition}[CE1 three-transverse return certificate]
\label{prop:app-ce1-direct-certificate}
Use the signed CE1 center variables
\[
0<R<1,
\qquad w=1-R,
\qquad E=\sqrt{1-Rw},
\]
\[
\eta=1-E,
\qquad P=E(1-E),
\]
and write
\[
 A=\alpha,
 \qquad D=\delta.
\]
Assume
\[
 A>0,
 \qquad D>0,
 \qquad A+wD<P,
 \qquad D+RA\ge P.
 \label{eq:app-ce1-signed-domain}
\]
Put
\[
 X=R-D,
 \qquad m=\frac AR,
 \qquad h_0=\frac{\eta+A+D}{2R}.
 \label{eq:app-ce1-derived-data}
\]
Let $T_0$ be the unique supercritical role and let
$T_1,\ldots,T_5$ be nonsupercritical $\Vdzero$ roles.  Assume
\[
 A_1\ge X,
\]
\[
 B_i+A_{i+1}\ge1\quad(1\le i\le4),
 \qquad B_5+A_0\ge1,
\]
and
\[
 C_4\ge1-A,\qquad C_3\ge1-m,\qquad C_2\ge1-D.
 \label{eq:app-ce1-transverse-demands}
\]
Then
\[
 \boxed{A_0>1-h_0.}
 \label{eq:app-ce1-return}
\]
\end{proposition}

\begin{figure}[!htbp]
\centering
\resizebox{.92\linewidth}{!}{%
\begin{tikzpicture}[
  x=1cm,y=1cm,>=Latex,
  state/.style={draw=black!38,rounded corners=1.5pt,fill=black!2,
    minimum height=.78cm,text width=2.35cm,align=center,font=\scriptsize},
  hard/.style={draw=RadialTeal,rounded corners=1.5pt,fill=RadialTeal!7,
    minimum height=.78cm,text width=2.35cm,align=center,font=\scriptsize},
  arr/.style={-{Stealth[length=1.8mm]},draw=black!60,line width=.65pt},
  harr/.style={-{Stealth[length=1.8mm]},draw=RadialTeal,line width=.85pt}
]
  \node[panellabel,anchor=west] at (-.1,2.55)
    {CE1 reverse path: from a deficient terminal to the final contradiction};

  \node[state] (a0) at (0,1.35) {$A_0\le1-h_0$};
  \node[state] (b5) at (3.05,1.35) {$B_5\ge h_0$};
  \node[state] (b4) at (6.10,1.35) {$B_4\ge h_0$};
  \node[hard]  (b3) at (9.15,1.35) {$B_3>L_1$};

  \node[hard]  (b2) at (0,-.55) {$B_2>L_2>\epsilon(D)$};
  \node[hard]  (a2) at (3.05,-.55) {$A_2\le\epsilon(D)$};
  \node[state] (b1) at (6.10,-.55) {$B_1>1-X$};
  \node[state,draw=DemandRed,fill=DemandRed!7]
    (contra) at (9.15,-.55)
    {$A_1\ge X$, $A_1+B_1\le1$\\contradiction};

  \draw[arr] (a0)--node[above=3pt,figlabel,fill=white,inner sep=.6pt] {companion edge} (b5);
  \draw[arr] (b5)--node[above=3pt,figlabel,fill=white,inner sep=.6pt] {Vd0 monotonicity} (b4);
  \draw[harr] (b4)--node[above=3pt,figlabel,fill=white,inner sep=.6pt] {selected $Q_+$ at $T_4$} (b3);
  \draw[harr] (b3.east)--++(.55,0)--++(0,-1.10)--
    node[right=2pt,figlabel,align=left,fill=white,inner sep=.6pt] {selected $Q_+$ at $T_3$\\and the scalar estimate}
    ++(-10.25,0)--(b2.west);
  \draw[harr] (b2)--node[above=3pt,figlabel,fill=white,inner sep=.6pt] {threshold at $T_2$} (a2);
  \draw[arr] (a2)--node[above=3pt,figlabel,fill=white,inner sep=.6pt] {coverage of $e_{1,2}$} (b1);
  \draw[arr,draw=DemandRed] (b1)--(contra);

  \node[figlabel,atlasbox,align=left,anchor=west] at (1.15,-1.68)
    {grey arrows: edge coverage or nonsupercriticality;\qquad
     teal arrows: the genuinely local high-radial steps};
  \node[figlabel,align=center] at (5.25,-2.33)
    {$L_1$ and $L_2$ are direct selected-arc lower bounds; no composed reach map is introduced.};
\end{tikzpicture}%
}
\caption{Logical structure of the CE1 direct reverse-path certificate.  The
proof starts from the negation of the desired terminal bound at $T_0$, moves
through the actual boundary reaches, uses two selected-arc estimates and one
high-radial threshold, and returns to $T_1$ with incompatible incoming and
outgoing demands.}
\label{fig:ce1-reverse-path}
\end{figure}


\paragraph{Solution.}

\begin{proof}[Proof of \zcref{prop:app-ce1-direct-certificate}]
The sign difference in \eqref{eq:app-ce1-signed-domain} gives
\[
 RD>wA,
\]
so the center exit on $r_3$ is $m=A/R$.  Moreover
\[
 0<A<P<\frac18,
 \qquad 0<D<\frac R2,
 \qquad 0<m<\frac w2<\frac12.
 \label{eq:app-ce1-elementary-bounds}
\]
We first record
\[
 \epsilon(A)<X,
 \qquad h_0>A.
 \label{eq:app-ce1-first-bounds}
\]
Indeed, $A+wD<P$ and $D=R-X$ give $A<wX-\eta$.  Convexity of
\[
 wX-\frac{X}{2(1+X)}
\]
on $R/2<X<R$ gives
\[
 A<\frac{X}{2(1+X)}.
\]
The low-root estimates preceding \zcref{lem:app-direct-threshold} yield
$\epsilon(A)<X$.  Also
\[
 2R(h_0-A)=\eta+D+(1-2R)A>0;
\]
when $R>1/2$, use $A<P=\eta E$ for the last sign.

Suppose, contrary to \eqref{eq:app-ce1-return}, that
\[
 A_0\le1-h_0.
 \label{eq:app-ce1-return-negation}
\]
The assumed final handoff gives $B_5\ge h_0$.
Nonsupercriticality of $T_5$ and the handoff $B_4+A_5\ge1$ imply
\[
 \boxed{B_4\ge B_5\ge h_0.}
 \label{eq:app-ce1-b4-lower}
\]

The assumed transverse radial reach at $T_4$ is at least $1-A$.  Apply the exact local
catalogue to the reflected boundary pair $(B_4,A_4)$.  If the local point is
not on the selected $Q_+$ component, the constant, $Q_-$, and linear
components together with \eqref{eq:app-ce1-first-bounds} give
\[
 1-A_4\ge1-\epsilon(A)>1-X.
\]
The center-free edges then carry this lower bound backward to
$B_1>1-X$, contradicting $A_1\ge X$ and
$A_1+B_1\le1$.

Hence the $T_4$ state is selected $Q_+$.  Lemma
\ref{lem:app-selected-chords} and \eqref{eq:app-ce1-b4-lower} give
\[
 B_3>L_1,
 \qquad
 L_1=(2-4A)h_0-(1-4A)A.
 \label{eq:app-ce1-l1}
\]
The selected-state condition also forces
\[
 \boxed{X>\frac12.}
 \label{eq:app-ce1-x-half}
\]
To see this, selectedness gives
\[
 h_0\le\epsilon(A)<\frac{2A}{1-2A}.
\]
If $X\le1/2$, the lower CE1 inequality gives
\[
 2Rh_0=\eta+A+D\ge w(R+A).
\]
Consequently
\[
 f_R(A)<0,
 \qquad
 f_R(z)=w(R+z)(1-2z)-4Rz.
\]
The quadratic is strictly concave and $f_R(0)>0$.  If $R\le1/2$, then
\[
 f_R(P)\ge\frac{1-E}{2}
 (6E^3-2E^2-5E+2)>0.
\]
If $R>1/2$, then $A<A_*=w/2-\eta$ and
\[
 f_R(A_*)=\frac{1-E}{2}(E+11R-3ER-5)>0.
\]
Concavity contradicts $f_R(A)<0$ in both cases, proving
\eqref{eq:app-ce1-x-half}.  Thus
\[
 m<\frac w2<\frac14.
 \label{eq:app-ce1-m-quarter}
\]

If $B_3\ge1/2$, then $B_3>1-X$ and the preceding immediate contradiction
applies.  Assume $B_3<1/2$.  At $T_3$ the reflected incoming reach is
$B_3>h_0>m$, while its assumed transverse radial reach is at least $1-m$.  If this state is
not selected $Q_+$, the exact local catalogue gives
\[
 B_2\ge1-B_3>\frac12>1-X,
\]
and the remaining center-free edge again contradicts $T_1$.

Thus both $T_4$ and $T_3$ are selected $Q_+$.  A second use of Lemma
\ref{lem:app-selected-chords} gives
\[
 B_2>(2-5m)B_3-(1-5m)m.
\]
Since $2-5m>0$, \eqref{eq:app-ce1-l1} yields
\[
 \boxed{
 B_2>L_2,
 \qquad
 L_2=(2-5m)L_1-(1-5m)m.
 }
 \label{eq:app-ce1-l2}
\]

We now prove the scalar estimate
\[
 \boxed{
 D<\frac1{10},
 \qquad
 L_2>2D+3D^2>\epsilon(D),
 \qquad
 \epsilon(D)<X.
 }
 \label{eq:app-ce1-scalar-estimate}
\]
The $T_4$ selected-state inequality and the bound
$\epsilon(A)\le2A+5A^2$ give
\[
 D\le D_h:=A(4R-1)+10RA^2-\eta.
 \label{eq:app-ce1-dh}
\]
If $D\ge1/10$, then $A+wD<P$ gives
\[
 A<\overline A=\frac{w(5R-1)}{10}.
\]
The right side of \eqref{eq:app-ce1-dh} is increasing in $A$, while
$\eta>Rw/2$, so
\[
 D<U(R):=\overline A(4R-1)+10R\overline A^2-\frac{Rw}{2}.
\]
Direct expansion gives
\[
 \frac1{10}-U(R)=-\frac R{10}P_4(R),
\]
where
\[
 P_4(R)=25R^4-60R^3+26R^2+22R-14.
\]
For $y=2R-1\in(0,1)$,
\[
 16P_4(R)
 =25y^4-20y^3-106y^2+124y-39
 \le-101y^2+124y-39<0.
\]
Hence $U(R)<1/10$, a contradiction.  This proves $D<1/10$.

Put
\[
 \Psi(D)=L_2-2D-3D^2,
 \qquad
 J(D)=L_2-\frac{23}{10}D.
\]
Since $D<1/10$,
\[
 \Psi(D)-J(D)=3D\left(\frac1{10}-D\right)>0.
\]
For fixed $R,A$, differentiation of \eqref{eq:app-ce1-l2} gives
\[
 J'(D)=\frac{S(R,A)}{10R^2},
\]
where
\[
 S(R,A)=100A^2-(40R+50)A+R(20-23R).
\]
The nonempty interval $P-RA\le D\le D_h$ implies
\[
 A>\frac{4Rw}{25R-4}=:A_*.
\]
On the relevant interval $\partial S/\partial A<-45$, and
\[
 S(R,A)<S(R,A_*)
 =-\frac{Rq_3(R)}{(25R-4)^2},
\]
where
\[
 q_3(R)=8775R^3-14260R^2+7928R-1120.
\]
For $y=2R-1$,
\[
 8q_3(R)=8775y^3-2195y^2+997y+3007>0.
\]
Thus $J$ decreases with $D$, so $J(D)\ge J(D_h)$.

At $D=D_h$, one has $h_0=2A+5A^2$.  If $L_{2,h}$ is the corresponding value,
then
\[
 L_{2,h}-A(1+4R)=\frac{A}{R^2}Q(R,A),
\]
where
\begin{align*}
Q(R,A)={}&100A^3R-40A^2R^2-30A^2R+12AR^2-15AR+5A\\
&-4R^3+5R^2-R.
\end{align*}
This polynomial is concave in $A$ on $0\le A\le Rw/2$.  Its endpoint values
are
\[
 Q(R,0)=Rw(4R-1)>0
\]
and
\[
 Q\left(R,\frac{Rw}{2}\right)
 =\frac{Rw}{2}
 (25R^5-30R^4+20R^3-3R^2-7R+3)>0.
\]
The last positivity follows after putting $y=2R-1$:
\[
32p(R)=25y^5+65y^4+90y^3+106y^2-35y+5>0.
\]
Hence $L_{2,h}>A(1+4R)$.

Finally $A<P=\eta E$, $A<Rw/2$, and $1/E>1+Rw/2$ give
\[
 \frac{D_h}{A}
 <C_*:=4R-2+5R^2w-\frac{Rw}{2}.
\]
Moreover
\[
 1+4R-\frac{23}{10}C_*
 =\frac{h_3(R)}{20},
 \qquad
 h_3(R):=230R^3-253R^2-81R+112.
\]
Here \eqref{eq:app-ce1-x-half} gives $R>1/2$, and therefore
$h_3''(R)=1380R-506\ge184$.  Strong convexity about
$R_0=20/23$, where
\[
 h_3(R_0)=\frac{788}{529},
 \qquad
 h_3'(R_0)=\frac{17}{23},
\]
gives the exact lower bound
\[
 h_3(R)
 \ge h_3(R_0)+h_3'(R_0)(R-R_0)+92(R-R_0)^2
 \ge\frac{289695}{194672}>0.
\]
Thus $J(D_h)>0$, so
\[
 L_2>2D+3D^2.
\]
To compare this value with the selected low root, substitute
$z=2D+3D^2$ directly into
\[
 p_D(z):=z^2-(1-D)z+(1-D)^2-(1-D)^4.
\]
The result is
\[
 p_D(2D+3D^2)=D^2(8D^2+19D-2)<0
 \qquad(0<D<1/10).
\]
Thus $2D+3D^2$ lies strictly between the two roots of $p_D$, and in
particular $\epsilon(D)<2D+3D^2$.  The argument proving
$\epsilon(A)<X$ also gives $\epsilon(D)<X$.  This proves
\eqref{eq:app-ce1-scalar-estimate}.

The assumed transverse radial reach at $T_2$ is at least $1-D$.  Since
$B_2>\epsilon(D)$, Lemma~\ref{lem:app-direct-threshold}, applied after
reflection, gives
\[
 A_2\le\epsilon(D).
\]
Coverage of $e_{1,2}$ therefore forces
\[
 B_1\ge1-A_2\ge1-\epsilon(D)>1-X,
\]
contradicting $A_1\ge X$ and $A_1+B_1\le1$.  Hence
\eqref{eq:app-ce1-return-negation} is impossible.
\end{proof}

\paragraph{Body consequence.}
This strict return estimate is the long CE1 supplier used by
\zcref{lem:app-ce1-transverse-terminal}, and hence proves the compact body
statement \zcref{prop:readable-k410-ce1}.
